\lyseis{2.2}{Ισοδύναμα κλάσματα}
\sol{1} $\frac{2}{3}=\frac{2\cdot 2}{3\cdot 2}=\frac{4}{6}$, $\frac{2}{3}=\frac{6}{9}$ και $\frac{2}{3}=\frac{20}{30}$.
\sol{2} α) $\frac{3}{4}=\frac{15}{20}$ (επί 5) \quad β) $\frac{5}{7}=\frac{35}{49}$ (επί 7) \quad γ) $\frac{18}{24}=\frac{3}{4}$ (διά 6) \quad δ) $\frac{5}{9}=\frac{20}{36}$ (επί 4).
\sol{3} α) Ναι: $4\cdot 15=60=6\cdot 10$. β) Όχι: $5\cdot 32=160$, $8\cdot 15=120$. γ) Ναι: $7\cdot 27=189=9\cdot 21$. δ) Ναι: $12\cdot 21=252=18\cdot 14$. ε) Όχι: $3\cdot 7=21$, $5\cdot 4=20$.
\sol{4} α) $\frac{18}{24}=\frac{3}{4}$ (ΜΚΔ 6) \quad β) $\frac{35}{42}=\frac{5}{6}$ (ΜΚΔ 7) \quad γ) $\frac{48}{72}=\frac{2}{3}$ (ΜΚΔ 24) \quad δ) $\frac{100}{250}=\frac{2}{5}$ (ΜΚΔ 50) \quad ε) $\frac{56}{98}=\frac{4}{7}$ (ΜΚΔ 14).
\sol{5} Ανάγωγα είναι τα $\frac{9}{16}$, $\frac{15}{28}$ και $\frac{25}{36}$. Τα υπόλοιπα: $\frac{12}{18}=\frac{2}{3}$, $\frac{21}{35}=\frac{3}{5}$, $\frac{17}{51}=\frac{1}{3}$.
\sol{6} α) $\frac{75}{100}$ \quad β) $\frac{35}{100}$ \quad γ) $\frac{40}{100}$ \quad δ) $\frac{18}{100}$ \quad ε) $\frac{52}{100}$.
\sol{7} α) $\frac{10}{12}$ \quad β) $\frac{9}{12}$ \quad γ) $\frac{8}{12}$ \quad δ) $\frac{6}{12}$.
\sol{8} α) $\frac{8}{12}$ και $\frac{9}{12}$ \quad β) $\frac{20}{24}$ και $\frac{21}{24}$ \quad γ) $\frac{8}{10}$, $\frac{5}{10}$ και $\frac{7}{10}$ \quad δ) ΕΚΠ(6, 8, 12) = 24, άρα $\frac{4}{24}$, $\frac{9}{24}$ και $\frac{10}{24}$.
\sol{9} α) Σωστό. β) Λάθος: $\frac{3}{5}=\frac{6}{10}$. γ) Λάθος: διαιρούμε και τους δύο όρους με τον ίδιο αριθμό. δ) Σωστό. ε) Λάθος: $\frac{6}{9}=\frac{2}{3}$ (ΜΚΔ 3). στ) Σωστό.
\sol{10} $\frac{6}{8}=\frac{3}{4}$, $\frac{9}{12}=\frac{3}{4}$, και $\frac{4}{6}=\frac{2}{3}$. Με κοινό παρονομαστή 12: Ρήνη $\frac{9}{12}$, Φανή $\frac{9}{12}$, Ηλίας $\frac{9}{12}$, Στέλιος $\frac{8}{12}$. Οι Ρήνη, Φανή και Ηλίας έφαγαν το ίδιο και περισσότερο από τον Στέλιο.
\sol{11} Έφαγε $\frac{18}{24}$, και με ΜΚΔ 6 γίνεται $\frac{3}{4}$ της σοκολάτας.
\sol{12} $\frac{45}{60}=\frac{3}{4}$ \quad $\frac{20}{60}=\frac{1}{3}$ \quad $\frac{40}{60}=\frac{2}{3}$ \quad $\frac{12}{60}=\frac{1}{5}$.
\sol{13} $\frac{250}{1000}=\frac{1}{4}$ \quad $\frac{750}{1000}=\frac{3}{4}$ \quad $\frac{400}{1000}=\frac{2}{5}$ \quad $\frac{125}{1000}=\frac{1}{8}$.
\sol{14} Χωράνε 2 όγδοα στο τέταρτο και 4 δέκατα έκτα: $\frac{1}{4}=\frac{2}{8}=\frac{4}{16}$.
\sol{15} Το $\frac{1}{2}$ είναι η μισή πίτσα, ενώ τα $\frac{4}{5}$ είναι σχεδόν ολόκληρη. Δεν είναι ίσα. Ο Στέλιος πρόσθεσε το 3 και στους δύο όρους, αλλά πρέπει να πολλαπλασιάσει: $\frac{1}{2}=\frac{1\cdot 3}{2\cdot 3}=\frac{3}{6}$.
\sol{16} Πολλαπλασιάζω τους όρους του $\frac{3}{5}$ με 1, 2, 3, \dots: $\frac{3}{5}$, $\frac{6}{10}$, $\frac{9}{15}$, $\frac{12}{20}$, $\frac{15}{25}$, $\frac{18}{30}$, $\frac{21}{35}$, $\frac{24}{40}$, $\frac{27}{45}$. Είναι 9 κλάσματα (το επόμενο, $\frac{30}{50}$, έχει παρονομαστή 50).
\sol{17} Το κλάσμα είναι $\frac{5k}{8k}$ και $5k+8k=13k=91$, άρα $k=7$ και το κλάσμα είναι $\frac{35}{56}$.
\sol{18} Το κλάσμα είναι $\frac{3k}{7k}$ και $7k-3k=4k=20$, άρα $k=5$ και το κλάσμα είναι $\frac{15}{35}$.
\sol{19} Αν τα 2 και 3 είχαν κοινό διαιρέτη $d>1$, τότε το $42\cdot d$ θα ήταν κοινός διαιρέτης των 84 και 126, και μεγαλύτερος από το 42. Αυτό είναι αδύνατο, γιατί το 42 είναι ο \emph{μέγιστος} κοινός διαιρέτης. Άρα $\MKD(2,\,3)=1$ και το κλάσμα είναι ανάγωγο.
\sol{20} Ισοδύναμα με το 0: $\frac{0}{3}$, $\frac{0}{7}$, $\frac{0}{12}$. Ισοδύναμα με το 1: $\frac{4}{4}$, $\frac{9}{9}$, $\frac{15}{15}$.
\sol{21} $\frac{2}{3}$, $\frac{4}{6}$, $\frac{6}{9}$, $\frac{8}{12}$: όλα είναι ίσα με $\frac{2}{3}$ (κάθε φορά πολλαπλασιάζονται οι όροι $2$ και $3$ με το ίδιο νούμερο).
\sol{22} Με ένα δίπλωμα: $\frac{1}{2}$ του φύλλου κάθε μέρος. Με δύο διπλώματα: $\frac{1}{4}$, και τα δύο μέρη μαζί $\frac{2}{4}$. Με τρία: $\frac{1}{8}$, και τέσσερα μέρη μαζί $\frac{4}{8}$. Άρα $\frac{1}{2}=\frac{2}{4}=\frac{4}{8}=\frac{8}{16}$.
